\cleardoublepage
\phantomsection
\pdfbookmark{Summary}{Summary}
\begingroup
\let\clearpage\relax
\let\cleardoublepage\relax
\let\cleardoublepage\relax

\chapter*{Summary}


Over the past few decades, the use of mobile devices has increased drastically, almost replacing desktop software and web browsers. This shift has been driven not only by the widespread adoption of smartphones, but also by the availability of integrated sensors such as GPS, gyroscopes, and Bluetooth, which allows a richer and more context-aware user experiences. By leveraging these technologies, Online Social Media platforms have been able to provide more advanced and sophisitcated features; consequently collecting larger amounts of metadata, interaction and daily engagement.

This thesis presents ``GeoMedia'', an open-source framework designed for sharing multimedia content — such as text, images, audio, and files — anchored to real-world geographic locations. By leveraging GPS data available on mobile devices, the application provides an interactive map through which users can explore and publish digital content, creating a seamless integration between the digital and physical worlds. Furthermore, the framework offers several features and cover various use cases, including restricting content visibility within customizable geographic areas, enabling time-sensitive content, applying user-based restrictions for content creation and visualization, and supporting sequential visibility between published posts. Furthermore, analyzes the motivations behind GeoMedia, its critical aspects and design considerations, describes its system architecture and implementation through a comparison of different technological paradigms. Finally, it explores potential applications and outlines possible future developments of the framework.

\vspace{1em}

An previous version of this framework has been presented at ``GoodIT '25: International Conference on Information Technology for Social Good'' (September 3 - 5, 2025 - Antwerp, Belgium) \cite{site:goodit}


%\vfill

%\selectlanguage{english}
%\pdfbookmark{Abstract}{Abstract}
%\chapter*{Abstract}

%\selectlanguage{italian}

\endgroup

\vfill
